\documentclass[11pt,a4paper]{article}

% ── Codificación, fuentes y lenguaje ───────────────────────────────────────────
\usepackage[utf8]{inputenc}
\usepackage[T1]{fontenc}
\usepackage[default,scale=0.95]{sourcesanspro}
\renewcommand{\familydefault}{\sfdefault}
\usepackage[spanish,es-noshorthands,es-tabla]{babel}

% ── Geometría y espaciado ─────────────────────────────────────────────────────
\usepackage[top=2.2cm, bottom=2.5cm, left=2.2cm, right=2.2cm, headheight=15pt]{geometry}
\usepackage{setspace}
\setstretch{1.18}
\usepackage{parskip}

% ── Matemáticas y unidades ────────────────────────────────────────────────────
\usepackage{amsmath}
\usepackage{amssymb}
\usepackage{siunitx}

% ── Circuitos ─────────────────────────────────────────────────────────────────
\usepackage[european, straightvoltages]{circuitikz}

% ── Tablas ────────────────────────────────────────────────────────────────────
\usepackage{booktabs}
\usepackage{tabularx}
\usepackage{array}
\usepackage{longtable}
\usepackage{colortbl}
\usepackage{enumitem}

% ── Colores, cajas e iconos ───────────────────────────────────────────────────
\usepackage[dvipsnames]{xcolor}
\usepackage{tcolorbox}
\tcbuselibrary{skins, breakable}
\usepackage{fontawesome5}
\usepackage{graphicx}

% ── Listados de código (dark-mode infográfico) ────────────────────────────────
%   Estilo base: fondo oscuro con números de línea. Los bloques lstlisting se
%   envuelven automáticamente en una caja tcolorbox redondeada.
\usepackage{listings}

% ── Secciones con barra de color (infografía) ────────────────────────────────
\usepackage{titlesec}
\titleformat{\section}
  {\normalfont\LARGE\bfseries\color{azulNoche}}
  {\colorbox{cyanNeon}{\color{fondoOscuro}\thesection}}{0.7em}{}
  [\vspace{2pt}\color{cyanNeon}\rule{\linewidth}{1.8pt}]
\titlespacing*{\section}{0pt}{24pt}{10pt}
\titleformat{\subsection}
  {\normalfont\large\bfseries\color{azulNoche}}
  {\textcolor{cyanNeon}{\thesubsection}}{0.7em}{}
  [\vspace{1pt}\color{grisLinea}\rule{\linewidth}{0.6pt}]
\titlespacing*{\subsection}{0pt}{14pt}{6pt}
\titleformat{\subsubsection}
  {\normalfont\normalsize\bfseries\color{azulNoche}}
  {\textcolor{cyanNeon}{\thesubsubsection}}{0.7em}{}
\titlespacing*{\subsubsection}{0pt}{10pt}{4pt}

% ── Cabeceras y pies de página ────────────────────────────────────────────────
\usepackage{fancyhdr}
\usepackage{lastpage}
\pagestyle{fancy}
\fancyhf{}
\renewcommand{\headrulewidth}{0pt}
\renewcommand{\footrulewidth}{0pt}
\fancyfoot[C]{%
  \begin{minipage}{\linewidth}
    \vspace{2pt}
    \color{cyanNeon}\rule{\linewidth}{1.2pt}\\[2pt]
    \centering\color{grisTexto}\footnotesize
    Página \thepage\ de \pageref{LastPage}
  \end{minipage}%
}

% ── Hiperenlaces ──────────────────────────────────────────────────────────────
\usepackage[colorlinks=true, linkcolor=azulNoche, urlcolor=cyanNeon]{hyperref}
\sloppy
\emergencystretch 3em

% ── Paleta de colores — tecnológico dark-mode ──────────────────────────────────
\definecolor{fondoOscuro}{HTML}{0D1B2A}%             Dark navy — fondo banda título
\definecolor{verdeTurquesa}{HTML}{004D40}%           Verde turquesa — bandas de marca
\definecolor{azulNoche}{HTML}{1B3A6B}%               Midnight blue — secciones, marcos
\definecolor{azulMedio}{HTML}{2A5298}%               Azul medio — variante de acento
\definecolor{cyanNeon}{HTML}{00C8E0}%                Cyan eléctrico — acentos primarios
\definecolor{verdeSignal}{HTML}{00B887}%             Verde señal — fortalezas, éxito
\definecolor{naranjaVivo}{HTML}{FF7043}%             Naranja vivo — mejoras, advertencias
\definecolor{rojoAlerta}{HTML}{E53935}%              Rojo alerta — errores
\definecolor{grisPapel}{HTML}{F4F6FA}%               Fondo tarjetas claras
\definecolor{grisLinea}{HTML}{C8D0E0}%               Bordes sutiles
\definecolor{grisTexto}{HTML}{6B7A99}%               Texto secundario
\definecolor{amarilloNota}{HTML}{FFB300}%            Notas / advertencias doradas
\definecolor{codigoFondo}{HTML}{1E2D3D}%             Fondo del bloque de código
\definecolor{codigoComentario}{HTML}{7A8FA6}%        Comentarios
\definecolor{codigoCadena}{HTML}{FF9D5C}%            Cadenas / strings
\definecolor{codigoKeyword}{HTML}{00C8E0}%           Palabras clave
\definecolor{codigoNumero}{HTML}{00E5A0}%            Números

% ── Banda de título: portada infográfica ──────────────────────────────────────
%   Diseño en 2 capas: banda de color (por defecto fondo oscuro) + franja de
%   acento cyan abajo. Uso: \bandaTitulo[color]{título}{subtítulo}.
%   El icono se puede cambiar con \renewcommand{\iconoBanda}{...} (FontAwesome).
\newcommand{\iconoBanda}{microchip}
\newcommand{\bandaTitulo}[3][fondoOscuro]{%
  \begin{tcolorbox}[
    enhanced,
    colback=#1, colframe=#1,
    boxrule=0pt, arc=8pt, outer arc=8pt,
    width=\linewidth,
    top=18pt, bottom=6pt, left=14pt, right=14pt,
    borderline south={3.5pt}{0pt}{cyanNeon},
  ]
    \begin{minipage}[c]{0.07\linewidth}
      \centering
      \textcolor{cyanNeon}{\fontsize{30}{30}\selectfont\faIcon{\iconoBanda}}
    \end{minipage}%
    \hspace{8pt}%
    \begin{minipage}[c]{0.88\linewidth}
      {\fontsize{20}{24}\selectfont\bfseries\color{white}#2}\par\vspace{5pt}%
      \textcolor{cyanNeon!80!white}{\small\faIcon{angle-right}~#3}%
    \end{minipage}
    \vspace{4pt}
  \end{tcolorbox}%
  \vspace{8pt}%
}

% ── Tarjetas de datos (infografía) ────────────────────────────────────────────
\tcbset{
  tarjetaDato/.style={
    enhanced, breakable,
    arc=6pt, outer arc=6pt,
    colback=grisPapel, colframe=azulNoche,
    boxrule=1pt,
    fonttitle=\bfseries\small, coltitle=white,
    attach boxed title to top left={yshift=-3mm, xshift=6mm},
    boxed title style={
      colback=azulNoche, colframe=azulNoche,
      arc=4pt, boxrule=0pt,
    },
    drop fuzzy shadow=azulNoche!30!white,
    top=8pt, bottom=8pt, left=10pt, right=10pt,
  },
  tarjetaIcono/.style={
    enhanced, breakable,
    arc=6pt, outer arc=6pt,
    colback=white, colframe=grisLinea, boxrule=0.8pt,
    fonttitle=\bfseries\small, coltitle=azulNoche,
    borderline west={3pt}{0pt}{cyanNeon},
    drop fuzzy shadow=azulNoche!25!white,
    top=6pt, bottom=6pt, left=10pt, right=10pt,
  },
  cajaTitulo/.style={
    enhanced, breakable,
    arc=6pt, outer arc=6pt,
    colback=azulNoche!10!grisPapel, colframe=azulNoche,
    boxrule=1pt,
    fonttitle=\bfseries, coltitle=white,
    attach boxed title to top left={yshift=-3mm, xshift=6mm},
    boxed title style={
      colback=azulNoche, colframe=azulNoche,
      arc=4pt, boxrule=0pt,
    },
    drop fuzzy shadow=azulNoche!25!white,
    top=8pt, bottom=8pt, left=10pt, right=10pt,
  },
  cajaContenido/.style={
    enhanced, breakable,
    arc=5pt, outer arc=5pt,
    colback=grisPapel, colframe=grisLinea,
    boxrule=0.8pt,
    fonttitle=\bfseries\small, coltitle=azulNoche,
    top=6pt, bottom=6pt, left=10pt, right=10pt,
  },
  cajaFortaleza/.style={
    enhanced, breakable,
    arc=6pt, outer arc=6pt,
    colback=verdeSignal!8!white, colframe=verdeSignal,
    boxrule=1pt,
    borderline west={4pt}{0pt}{verdeSignal},
    fonttitle=\bfseries\small, coltitle=verdeSignal!70!black,
    drop fuzzy shadow=verdeSignal!20!white,
    top=6pt, bottom=6pt, left=10pt, right=10pt,
  },
  cajaMejora/.style={
    enhanced, breakable,
    arc=6pt, outer arc=6pt,
    colback=naranjaVivo!8!white, colframe=naranjaVivo,
    boxrule=1pt,
    borderline west={4pt}{0pt}{naranjaVivo},
    fonttitle=\bfseries\small, coltitle=naranjaVivo!80!black,
    drop fuzzy shadow=naranjaVivo!20!white,
    top=6pt, bottom=6pt, left=10pt, right=10pt,
  },
  cajaRecomendacion/.style={
    enhanced, breakable,
    arc=6pt, outer arc=6pt,
    colback=cyanNeon!6!white, colframe=cyanNeon!60!azulNoche,
    boxrule=1pt,
    borderline west={4pt}{0pt}{cyanNeon},
    fonttitle=\bfseries\small, coltitle=azulNoche,
    drop fuzzy shadow=azulNoche!20!white,
    top=6pt, bottom=6pt, left=10pt, right=10pt,
  },
}

% ── Ícono + texto (encabezados de sección y elementos) ────────────────────────
\newcommand{\iconotexto}[2]{\textcolor{cyanNeon}{\faIcon{#1}}~#2}

% ── Listados de código: estilo dark + auto-envoltura ─────────────────────────
\lstdefinestyle{estiloCodigo}{
    backgroundcolor=\color{codigoFondo},
    basicstyle=\ttfamily\small\color{white},
    commentstyle=\color{codigoComentario}\itshape,
    stringstyle=\color{codigoCadena},
    keywordstyle=\color{codigoKeyword}\bfseries,
    numberstyle=\tiny\color{codigoComentario},
    numbers=left,
    numbersep=10pt,
    showstringspaces=false,
    breaklines=true,
    breakatwhitespace=false,
    tabsize=4,
    frame=none,
    captionpos=b,
    aboveskip=6pt,
    belowskip=4pt,
    xleftmargin=14pt,
    framexleftmargin=14pt,
    literate=
      {á}{{\'a}}1 {é}{{\'e}}1 {í}{{\'i}}1 {ó}{{\'o}}1 {ú}{{\'u}}1
      {Á}{{\'A}}1 {É}{{\'E}}1 {Í}{{\'I}}1 {Ó}{{\'O}}1 {Ú}{{\'U}}1
      {ñ}{{\~n}}1 {Ñ}{{\~N}}1 {ü}{{\"u}}1 {Ü}{{\"U}}1
      {¿}{{?`}}1 {¡}{{!`}}1
}
\lstdefinestyle{estiloPython}{
    style=estiloCodigo,
    language=Python,
}
\lstset{style=estiloCodigo}
\BeforeBeginEnvironment{lstlisting}{%
  \begin{tcolorbox}[
    enhanced, breakable,
    arc=6pt, outer arc=6pt,
    colback=codigoFondo, colframe=azulNoche,
    boxrule=1pt,
    drop fuzzy shadow=azulNoche!25!white,
    top=2pt, bottom=2pt, left=0pt, right=0pt,
  ]%
}
\AfterEndEnvironment{lstlisting}{\end{tcolorbox}}

% ── Cajas de contenido temáticas ──────────────────────────────────────────────
\newtcolorbox{cajaRecuerda}{
    enhanced, breakable,
    arc=6pt, outer arc=6pt,
    colback=azulNoche!10!grisPapel,
    colframe=azulNoche, boxrule=1pt,
    borderline west={4pt}{0pt}{cyanNeon},
    fonttitle=\bfseries\small\color{white},
    title={\faIcon{info-circle}~Recuerda},
    attach boxed title to top left={yshift=-3mm, xshift=6mm},
    boxed title style={colback=azulNoche, arc=4pt, boxrule=0pt},
    drop fuzzy shadow=azulNoche!25!white,
    left=10pt, right=10pt, top=8pt, bottom=8pt,
}

\newtcolorbox{cajaConcepto}{
    enhanced, breakable,
    arc=6pt, outer arc=6pt,
    colback=verdeSignal!8!white,
    colframe=verdeSignal, boxrule=1pt,
    borderline west={4pt}{0pt}{verdeSignal},
    fonttitle=\bfseries\small\color{white},
    title={\faIcon{lightbulb}~Concepto clave},
    attach boxed title to top left={yshift=-3mm, xshift=6mm},
    boxed title style={colback=verdeSignal!80!black, arc=4pt, boxrule=0pt},
    drop fuzzy shadow=verdeSignal!20!white,
    left=10pt, right=10pt, top=8pt, bottom=8pt,
}

\newtcolorbox{cajaEjemplo}{
    enhanced, breakable,
    arc=6pt, outer arc=6pt,
    colback=cyanNeon!5!white,
    colframe=cyanNeon!70!azulNoche, boxrule=1pt,
    borderline west={4pt}{0pt}{cyanNeon},
    fonttitle=\bfseries\small\color{fondoOscuro},
    title={\faIcon{code}~Ejemplo},
    attach boxed title to top left={yshift=-3mm, xshift=6mm},
    boxed title style={colback=cyanNeon, arc=4pt, boxrule=0pt},
    drop fuzzy shadow=azulNoche!20!white,
    left=10pt, right=10pt, top=8pt, bottom=8pt,
}

\newtcolorbox{cajaEjercicio}{
    enhanced, breakable,
    arc=6pt, outer arc=6pt,
    colback=azulMedio!7!white,
    colframe=azulMedio, boxrule=1pt,
    borderline west={4pt}{0pt}{azulMedio},
    fonttitle=\bfseries\small\color{white},
    title={\faIcon{pencil-alt}~Ejercicio},
    attach boxed title to top left={yshift=-3mm, xshift=6mm},
    boxed title style={colback=azulMedio, arc=4pt, boxrule=0pt},
    drop fuzzy shadow=azulNoche!20!white,
    left=10pt, right=10pt, top=8pt, bottom=8pt,
}

\newtcolorbox{cajaReto}{
    enhanced, breakable,
    arc=6pt, outer arc=6pt,
    colback=naranjaVivo!6!white,
    colframe=naranjaVivo, boxrule=1pt,
    borderline west={4pt}{0pt}{naranjaVivo},
    fonttitle=\bfseries\small\color{white},
    title={\faIcon{fire}~Reto},
    attach boxed title to top left={yshift=-3mm, xshift=6mm},
    boxed title style={colback=naranjaVivo, arc=4pt, boxrule=0pt},
    drop fuzzy shadow=naranjaVivo!20!white,
    left=10pt, right=10pt, top=8pt, bottom=8pt,
}

% ── Indicadores de dificultad ─────────────────────────────────────────────────
\newcommand{\facil}{\textcolor{verdeSignal}{\faIcon{circle}\,\textbf{Fácil}}}
\newcommand{\intermedio}{\textcolor{naranjaVivo}{\faIcon{circle}\,\textbf{Intermedio}}}
\newcommand{\dificil}{\textcolor{rojoAlerta}{\faIcon{circle}\,\textbf{Difícil}}}

% ─────────────────────────────────────────────────────────────────────────────
\renewcommand{\iconoBanda}{atom}
\fancyhead[L]{\footnotesize\color{grisTexto}\faIcon{atom}~LABORATORIO\_II\_FISICA}
\fancyhead[R]{\footnotesize\color{grisTexto}Práctica 01}
\hypersetup{pdftitle={Laboratorio II de Física},pdfauthor={Prof. César Rodríguez}
  pdfsubject={LABORATORIO\_II\_FISICA}}

\begin{document}
\bandaTitulo[fondoOscuro]{Práctica 01 --- Característica corriente-tensión de un resistor}{Semana 1 --- Fase 1: Instrumentación y leyes básicas de los circuitos}

\begin{center}\small\color{grisTexto}Prof. César Rodríguez \textbar{} Universidad de Oriente, Núcleo de Sucre, Departamento de Física \textbar{} 2 de noviembre de 2026\end{center}

\begin{tabularx}{\linewidth}{@{}>{\raggedright\arraybackslash}p{4.6cm}X@{}}
\toprule
\textbf{Duración} & 3 horas de laboratorio presencial más 2 horas de análisis \\
\textbf{Modalidad} & 4 grupos con roles rotativos \\
\textbf{Entrega} & GitHub del grupo, carpeta semana\_01, antes del corte semanal \\
\textbf{Instrumento} & Práctica formativa con informe \\
\bottomrule
\end{tabularx}

\section{\iconotexto{bullseye}{Objetivo}}

Aprender a usar la fuente de poder, los multimedidores y el tablero de montaje, estimar la incertidumbre de cada lectura y decidir, a partir de la característica $I$--$V$, si un elemento es óhmico o no lineal.

\section{\iconotexto{clipboard-list}{Contenidos que se aplican}}

\begin{itemize}[leftmargin=1.4em,itemsep=2pt,topsep=3pt]
  \item Magnitudes eléctricas y unidades del SI; prefijos y notación científica.
  \item Resolución, exactitud y clase de un instrumento; dígitos de un multímetro.
  \item Incertidumbre tipo A y tipo B; expresión del resultado con cifras significativas.
  \item Conexión del voltímetro en paralelo y del amperímetro en serie; efecto de carga.
  \item Ley de Ohm y criterio experimental de ohmicidad.
  \item Ajuste por mínimos cuadrados, con pendiente, ordenada y sus incertidumbres.
  \item Normas de seguridad eléctrica en la mesa de laboratorio.
\end{itemize}

\section{\iconotexto{book-open}{Fundamento teórico}}

Un resistor es óhmico cuando la corriente que lo atraviesa es proporcional
a la tensión entre sus bornes:
\[
  V = R\,I \qquad\Longrightarrow\qquad R = \frac{V}{I}
\]
La resistencia no se estima como el promedio de los cocientes $V_k/I_k$,
sino como la pendiente del ajuste lineal de $V$ frente a $I$, que pondera
todos los puntos y reduce el ruido de lectura:
\[
  R = \frac{n\sum_k V_k I_k - \sum_k V_k \sum_k I_k}
           {n\sum_k I_k^2 - \left(\sum_k I_k\right)^2}
\]
La incertidumbre de la lectura de un multímetro digital combina un
porcentaje de la lectura con un número de dígitos del último lugar:
\[
  u(V) = a \cdot V + b \cdot \mathrm{d\acute{i}g}
\]
con $a$ expresada en tanto por uno y $b$ en dígitos del último lugar.
y, al propagar a una magnitud indirecta $R = V/I$,
\[
  \frac{u_R}{R} = \sqrt{\left(\frac{u_V}{V}\right)^2
                      + \left(\frac{u_I}{I}\right)^2}
\]
La dispersión de los puntos alrededor del ajuste se mide con los
residuales $e_k = V_k - R I_k$; si los residuales tienen signo
sistemático en función de la corriente, el elemento no es óhmico o hay un
error sistemático, como calentamiento del resistor o efecto de carga del
amperímetro.

\begin{center}
\begin{circuitikz}
  \draw (0,0) to[V, l=$E$] (0,3);
  \draw (0,3) to[ammeter, l=$I$] (3,3);
  \draw (3,3) to[R, l=$R$] (3,0);
  \draw (3,0) -- (0,0);
  \draw (3,3) -- (4.6,3) to[voltmeter, l=$V_R$] (4.6,0) -- (3,0);
\end{circuitikz}
\par\smallskip{\small\itshape Montaje de la práctica 1, con amperímetro en serie y voltímetro en paralelo con el resistor}
\end{center}

\section{\iconotexto{box-open}{Materiales y equipos}}

\begin{itemize}[leftmargin=1.4em,itemsep=2pt,topsep=3pt]
  \item Una fuente de poder de 0 a 30 V y 6 A
  \item Dos multimedidores Keithley (179, 175 o 175-A)
  \item Un diodo semiconductor rectificador
  \item Una resistencia variable o de carbón
  \item Tablero de montaje y cables de conexión
  \item Manual de instrucciones de los Keithley 179 y 175
\end{itemize}

\section{\iconotexto{flask}{Escenario}}

Cada grupo recibe un resistor de carbón, un diodo semiconductor y dos multimedidores Keithley. Debe obtener la característica $I$--$V$ de los dos dispositivos sobre el tablero, decidir cuál cumple la ley de Ohm y reportar la resistencia del resistor con la incertidumbre del ajuste, no con el promedio de los cocientes.

\section{\iconotexto{list-ol}{Procedimiento}}

\begin{enumerate}[leftmargin=1.4em,itemsep=2pt,topsep=3pt]
  \item Leer, en el manual de los Keithley, cómo se selecciona cada función, cada rango y la resolución en voltímetro, amperímetro y óhmetro.
  \item Identificar el valor nominal del resistor por código de colores y anotarlo con su tolerancia.
  \item Montar el circuito del resistor, sin energizar, y hacerlo revisar por el profesor.
  \item Conectar el voltímetro en paralelo con el resistor y el amperímetro en serie, ambos en el rango más alto, y bajar el rango hasta resolver la lectura.
  \item Para un valor fijo de R, variar la salida de la fuente y obtener al menos ocho pares $V$, $I$, anotando rango y resolución en cada punto.
  \item Sustituir el resistor por el diodo, en polarización directa, y obtener un mínimo de veinte pares $V$, $I$ variando la fuente o la resistencia.
  \item Invertir la polarización del diodo y comprobar que la corriente es prácticamente nula dentro del rango del amperímetro.
  \item Cortar la alimentación, medir el resistor con el óhmetro y comparar con el valor nominal y con el del ajuste.
  \item Verificar el efecto de carga repitiendo tres puntos con el voltímetro antes del amperímetro y anotando la diferencia.
\end{enumerate}

\section{\iconotexto{table}{Tabla de datos}}

\begin{center}
\begin{minipage}{\linewidth}
\centering
{\small
\textbf{Registro de la característica corriente-tensión de los dos elementos}\par\smallskip
\begin{tabular}{@{}cccccc@{}}
\toprule
Elemento & Ensayo & $V$ (V) & $I$ (mA) & $u_V$ (V) & $u_I$ (mA) \\
\midrule
Resistor & 1 & 1,00 & $\ldots$ & $\ldots$ & $\ldots$ \\
Resistor & 2 & 2,00 & $\ldots$ & $\ldots$ & $\ldots$ \\
Resistor & 3 & 4,00 & $\ldots$ & $\ldots$ & $\ldots$ \\
Resistor & 4 & 10,00 & $\ldots$ & $\ldots$ & $\ldots$ \\
Diodo & 1 & 0,20 & $\ldots$ & $\ldots$ & $\ldots$ \\
Diodo & 2 & 0,40 & $\ldots$ & $\ldots$ & $\ldots$ \\
Diodo & 3 & 0,55 & $\ldots$ & $\ldots$ & $\ldots$ \\
Diodo & 4 & 0,65 & $\ldots$ & $\ldots$ & $\ldots$ \\
\bottomrule
\end{tabular}
}\par\smallskip{\footnotesize Mínimo veinte filas para el diodo; anotar la resolución de cada instrumento y entregar los datos crudos en CSV.}
\end{minipage}
\end{center}

\section{\iconotexto{calculator}{Análisis y tratamiento de datos}}

\begin{enumerate}[leftmargin=1.4em,itemsep=2pt,topsep=3pt]
  \item Graficar $I$ frente a $V$ para cada elemento, con barras de incertidumbre y ejes rotulados con unidades.
  \item Ajustar por mínimos cuadrados la rama lineal del resistor y reportar la pendiente, su incertidumbre y el coeficiente de determinación.
  \item Calcular los residuales del resistor y decidir si son aleatorios o sistemáticos; justificar con la gráfica.
  \item Comparar la resistencia del ajuste con el valor nominal, con el óhmetro y con el código de colores.
  \item Describir la forma de la curva del diodo en directa y en inversa, y concluir si cumple o no la ley de Ohm.
  \item Estimar la potencia máxima disipada por el resistor y comprobar que no excede su potencia nominal.
  \item Cuantificar el efecto de carga observado en el último paso y atribuirlo a la resistencia interna del instrumento.
\end{enumerate}

\section{\iconotexto{box-open}{Entregable}}

Tabla de datos crudos en CSV, gráfica de la característica con ajuste y banda de incertidumbre, script que reproduce la figura, y valor de la resistencia con su incertidumbre comparado contra el nominal.

\section{\iconotexto{ruler-combined}{Rúbrica de corrección}}

\begin{tabularx}{\linewidth}{@{}X r@{}}
\toprule
Criterio & Peso\\
\midrule
Fundamento teórico y modelo aplicado & 1,0 \\
Montaje, seguridad y manejo del instrumental & 1,5 \\
Toma de datos completa y trazable & 2,0 \\
Presupuesto de incertidumbres & 2,0 \\
Análisis, ajuste y gráficas & 2,0 \\
Conclusiones contrastadas con el modelo & 1,0 \\
Presentación y estructura del informe & 0,5 \\
\midrule
\textbf{Total} & \textbf{10,0}\\
\bottomrule
\end{tabularx}

\section{\iconotexto{code}{Contrato del script}}

\begin{cajaConcepto}Todo script entregado imprime \textbf{un solo objeto JSON} por la salida estándar y devuelve un diccionario con estas claves. \texttt{status} vale \texttt{ok} o \texttt{error}; si es \texttt{error}, \texttt{errores} explica el motivo y el programa sale con código distinto de cero.
\begin{lstlisting}[language=Python,style=estiloPython]
{
  "status": "ok",
  "practica": 1,
  "magnitud": "resistencia",
  "valor": 99.4,
  "unidad": "ohm",
  "incertidumbre": 0.6,
  "metodo": "ajuste_lineal_minimos_cuadrados",
  "n_puntos": 8,
  "residual_max": 0.012,
  "fuente_datos": "semana_01/datos/iv_resistor_220.csv",
  "errores": []
}

\end{lstlisting}
\end{cajaConcepto}

\begin{tcolorbox}[cajaMejora]\textbf{Antes de empezar:} Antes de energizar, contrasta el montaje con la figura y verifica rangos y polaridad de cada instrumento. Si un instrumento se calienta, huele o marca fuera de escala, corta la alimentación y avisa al docente: el montaje se revisa antes de continuar.
\end{tcolorbox}
\end{document}
